% Archive record 34, line 1.
\hypertarget{part-3-34-0001-iii-b-tools-memory-planning-reflection}{%
\section{III-B --- Tools, Memory, Planning \& Reflection}\label{part-3-34-0001-iii-b-tools-memory-planning-reflection}}

上一阶段 III-A，我们将 AI Agent 抽象为：
\[A=(\mathcal L,P,M,T,C).\]

其中 \(\mathcal L\) 是底层语言模型，\(P\) 是 Prompt，\(M\) 是 Memory，\(T\) 是 Tools，\(C\) 是 Control Logic。

本阶段，我们将进一步研究 Agent 如何利用这些组件完成复杂任务，以及它们如何影响 RSI 的泛化能力。

\begin{readertable}
\begin{longtable}[]{@{}
  >{\raggedright\arraybackslash}p{(\columnwidth - 6\tabcolsep) * \real{0.18}}
  >{\raggedright\arraybackslash}p{(\columnwidth - 6\tabcolsep) * \real{0.27}}
  >{\raggedright\arraybackslash}p{(\columnwidth - 6\tabcolsep) * \real{0.37}}
  >{\raggedright\arraybackslash}p{(\columnwidth - 6\tabcolsep) * \real{0.15}}@{}}
\toprule
Component & Core Function & 与 RSI 的联系 & Priority \\
\midrule
\endhead
Tools & 执行外部操作 & 扩大可搜索的方法空间 & P1 \\
Memory & 保存与利用历史信息 & 累积经验，也可能累积评估偏差 & P0 \\
Planning & 组织未来行动 & 决定搜索方向和资源分配 & P1 \\
Reflection & 根据反馈修正行为 & 形成 Self-Improvement 的基础 & P0 \\
\bottomrule
\end{longtable}
\end{readertable}

其中 \textbf{Memory 和 Reflection} 是我们本阶段的重点，因为这两项机制与 II-D 中的 Adaptive Data Analysis 直接相关。

\hypertarget{part-3-34-0001-agent-ux7684ux57faux672cux8ba4ux77e5ux4e0eux6267ux884cux5faaux73af}{%
\subsection{1. Agent 的基本认知与执行循环}\label{part-3-34-0001-agent-ux7684ux57faux672cux8ba4ux77e5ux4e0eux6267ux884cux5faaux73af}}

一个 Financial Research Agent 在完成任务时，可能需要经历以下过程。

\bookgraphic{0.342}{part-3-34-0001-001-01}{上方是 Agent 的任务执行循环（Inner Loop）；下方的持久性修改可能成为 Outer Improvement Loop 的一部分。并非每次 Reflection 都会产生新的 Agent 版本。}

这四项能力分别回答四个问题：

\begin{itemize}
\tightlist
\item
  \textbf{Planning：} 下一步应该做什么？
\item
  \textbf{Tools：} 怎样真正执行这一步？
\item
  \textbf{Memory：} 过去已经做过什么？
\item
  \textbf{Reflection：} 应该如何根据结果改变未来的行为？
\end{itemize}

单独拥有其中某项能力，并不意味着 Agent 具备 Recursive Self-Improvement。真正重要的是这些能力能否形成一个持续的、基于反馈的改进闭环。

\hypertarget{part-3-34-0001-tools-agent-ux5982ux4f55ux4e0eux5916ux90e8ux4e16ux754cux4ea4ux4e92}{%
\subsection{2. Tools --- Agent 如何与外部世界交互？}\label{part-3-34-0001-tools-agent-ux5982ux4f55ux4e0eux5916ux90e8ux4e16ux754cux4ea4ux4e92}}

\hypertarget{part-3-34-0001-tool-calling-ux7684ux57faux672cux539fux7406}{%
\subsubsection{2.1 Tool Calling 的基本原理}\label{part-3-34-0001-tool-calling-ux7684ux57faux672cux539fux7406}}

LLM 本身负责生成 Token，但 Token 可以被 Harness 解释为结构化操作。

例如，Agent 希望检验某个金融特征的预测能力，于是生成如下 Tool Call：

\begin{verbatim}
{
    "tool": "run_experiment",
    "arguments": {
        "feature": "momentum",
        "lookback": 20,
        "model": "ridge_regression"
    }
}
\end{verbatim}

Harness 接收到该请求后，检查参数与权限，执行相应的 Python 函数，并返回实验结果：

\begin{verbatim}
{
    "status": "success",
    "is_score": 0.638,
    "runtime_seconds": 2.41
}
\end{verbatim}

LLM 根据这些结果决定下一步操作。

因此，Tool Calling 的本质是：
\[\text{Language Generation}
\longrightarrow
\text{Structured Action}
\longrightarrow
\text{Program Execution}
\longrightarrow
\text{Observation}.\]

上述接口是教学示例，并非课题组已有的实际 API。

\hypertarget{part-3-34-0001-research-agent-ux7684ux5178ux578bux5de5ux5177}{%
\subsubsection{2.2 Research Agent 的典型工具}\label{part-3-34-0001-research-agent-ux7684ux5178ux578bux5de5ux5177}}

\begin{readertable}
\begin{longtable}[]{@{}
  >{\raggedright\arraybackslash}p{(\columnwidth - 4\tabcolsep) * \real{0.31}}
  >{\raggedright\arraybackslash}p{(\columnwidth - 4\tabcolsep) * \real{0.29}}
  >{\raggedright\arraybackslash}p{(\columnwidth - 4\tabcolsep) * \real{0.36}}@{}}
\toprule
Tool Category & Function & Financial Example \\
\midrule
\endhead
Data Access & 获取和读取数据 & 读取历史价格 \\
Code Execution & 执行程序 & 运行 Regression \\
File Operations & 创建及修改代码 & 修改特征生成函数 \\
Evaluation & 评价候选方案 & 计算 Validation MSE \\
Version Control & 保存及恢复版本 & Git Commit、Rollback \\
Retrieval & 检索外部知识 & 查询研究论文 \\
Resource Monitor & 记录资源消耗 & Runtime、Token Usage \\
\bottomrule
\end{longtable}
\end{readertable}

\hypertarget{part-3-34-0001-tools-ux4e0e-statistical-complexity}{%
\subsubsection{2.3 Tools 与 Statistical Complexity}\label{part-3-34-0001-tools-ux4e0e-statistical-complexity}}

新增工具不仅意味着 Agent 获得了新的执行能力，也可能改变它所能够探索的假设空间。

例如：

\begin{readertable}
\begin{longtable}[]{@{}
  >{\raggedright\arraybackslash}p{(\columnwidth - 4\tabcolsep) * \real{0.38}}
  >{\raggedright\arraybackslash}p{(\columnwidth - 4\tabcolsep) * \real{0.22}}
  >{\raggedright\arraybackslash}p{(\columnwidth - 4\tabcolsep) * \real{0.36}}@{}}
\toprule
& Agent \(A_0\) & Agent \(A_1\) \\
\midrule
\endhead
Regression Tool & 支持 & 支持 \\
Feature Generation & 手动生成 & 自动批量生成 \\
Model Evaluation & 单个模型 & 并行比较大量模型 \\
Search Space & 相对有限 & 可能显著扩大 \\
\bottomrule
\end{longtable}
\end{readertable}

如果新工具能够访问更多候选方法，那么可能有：
\[\mathcal H_0\subsetneq\mathcal H_1.\]

其中 \(\mathcal H_k\) 表示第 \(k\) 个版本在给定约束下能够搜索的假设空间。

结合 II-B：
\[\epsilon=
\sqrt{\frac{\log(2|\mathcal H|/\alpha)}{2n}}.\]

在固定有限假设空间、有界损失及 IID 等条件下，扩大搜索空间会使这一泛化误差上界变得更加宽松。 因此：

\textbf{Tool Expansion 既可能是 Capability Improvement，也可能是 Overfitting Amplification。}

这意味着，Revision Classification 中的 Infrastructure Improvement 不能一律被视为与统计性能无关的修改。

例如，某个 Revision 将实验执行速度提高十倍。在相同时间预算下，Agent 可能因此测试十倍数量的候选策略，间接改变 Model Selection Pressure。

\hypertarget{part-3-34-0001-memory-agent-ux5982ux4f55ux79efux7d2fux7ecfux9a8c}{%
\subsection{3. Memory --- Agent 如何积累经验？}\label{part-3-34-0001-memory-agent-ux5982ux4f55ux79efux7d2fux7ecfux9a8c}}

这是本阶段第一个重点。

\hypertarget{part-3-34-0001-memory-ux4e0dux7b49ux4e8e-context-window}{%
\subsubsection{3.1 Memory 不等于 Context Window}\label{part-3-34-0001-memory-ux4e0dux7b49ux4e8e-context-window}}

Context Window 指单次 LLM 调用可以处理的上下文范围，而 Memory 是整个系统保存、组织和提取历史信息的机制。

一个 Agent 可以拥有很大的外部 Memory，却在每一次调用 LLM 时只检索其中极少量相关信息。

\bookgraphic{0.463}{part-3-34-0001-002-01}{Memory 可以长期存在于数据库、文件或版本仓库中，Context Window 只是当前调用所能看到的信息集合。}

这也涉及一个比较重要的工程概念：\textbf{Context Engineering}。

Prompt Engineering 主要关注如何编写指令，而 Context Engineering 还涉及决定哪些数据、历史记录、工具说明及中间状态应该进入当前上下文。

\hypertarget{part-3-34-0001-four-types-of-agent-memory}{%
\subsubsection{3.2 Four Types of Agent Memory}\label{part-3-34-0001-four-types-of-agent-memory}}

我们可以采用一种常见但并非统一标准的功能性分类。

\begin{readertable}
\begin{longtable}[]{@{}
  >{\raggedright\arraybackslash}p{(\columnwidth - 4\tabcolsep) * \real{0.27}}
  >{\raggedright\arraybackslash}p{(\columnwidth - 4\tabcolsep) * \real{0.32}}
  >{\raggedright\arraybackslash}p{(\columnwidth - 4\tabcolsep) * \real{0.37}}@{}}
\toprule
Memory Type & 含义 & Research Example \\
\midrule
\endhead
Working Memory & 当前任务的临时状态 & 正在执行的实验 \\
Episodic Memory & 过去发生的具体事件 & Round 5 的实验失败原因 \\
Semantic Memory & 从经验中提炼的知识 & 正则化可以限制模型复杂度 \\
Procedural Memory & 完成任务的方法和程序 & 自动化模型搜索流程 \\
\bottomrule
\end{longtable}
\end{readertable}

对于 RSI，后面三种尤其值得研究。

例如，同一个实验结果可以在 Memory 中以三种形式存在。

\textbf{Episodic Memory}

\begin{verbatim}
Round 5:
Increasing the number of features from 10 to 50
improved IS performance by 0.03.
\end{verbatim}

\textbf{Semantic Memory}

\begin{verbatim}
Increasing feature diversity may improve
financial prediction performance.
\end{verbatim}

\textbf{Procedural Memory}

\begin{verbatim}
def generate_features(data):
    return generate_50_features(data)
\end{verbatim}

这三个阶段形成了一种潜在的知识转化过程：
\[\text{Experience}
\longrightarrow
\text{Generalization of Experience}
\longrightarrow
\text{Persistent Procedure}.\]

但注意，中间的 Generalization of Experience 是 Agent 对实验现象的归纳，不等于我们此前讨论的统计泛化保证。

这正是危险所在。

如果原来的 IS Improvement 只是随机波动，那么 Agent 可能将一次偶然的成功转化为一个错误的长期研究原则，甚至进一步将其写入自身代码。

\hypertarget{part-3-34-0001-mathematical-representation-of-memory-updating}{%
\subsubsection{3.3 Mathematical Representation of Memory Updating}\label{part-3-34-0001-mathematical-representation-of-memory-updating}}

沿用 Part I 中的历史记录 \(H_k\)。

设 \(F_k\) 为第 \(k\) 轮获得的实验反馈，\(M_k\) 为该版本所使用的持久性 Memory。

那么可以定义：
\[\boxed{
M_{k+1}
=
U_M(M_k,F_k,H_k)
}\]

其中 \(U_M\) 表示 Memory Update Mechanism。

下一轮 Agent 的行为则受到 \(M_{k+1}\) 的影响。

\[A_{k+1}
=
(\mathcal L,P_{k+1},M_{k+1},T_{k+1},C_{k+1}).\]

从计算机科学角度来看，这是正常的状态持久化与知识积累。

但从统计学角度来看，这同时意味着：
\[D_{\mathrm{IS}}
\longrightarrow F_k
\longrightarrow M_{k+1}
\longrightarrow A_{k+1}.\]

这与 II-D 中的 Adaptive Information Flow 完全对应。

Memory 成为了连接历史评估数据与未来 Agent 行为的一种信息通道。

\hypertarget{part-3-34-0001-memory-ux7684ux53ccux91cdux4f5cux7528}{%
\subsubsection{3.4 Memory 的双重作用}\label{part-3-34-0001-memory-ux7684ux53ccux91cdux4f5cux7528}}

\textbf{Beneficial Memory}

Reliable experimental evidence

↓

Transferable knowledge

↓

Better future decisions

减少重复错误，积累可迁移的研究经验。

\textbf{Harmful Memory}

Noisy IS improvement

↓

Spurious research principle

↓

Repeated biased decisions

将偶然规律转化为持续影响搜索行为的错误经验。

这里可以提出一个研究假设：

\textbf{RSI 的后期性能退化，可能不仅来自单次错误修改，也来自错误经验在 Memory 中持续积累。}

我们暂时将这种候选机制称为 \textbf{Memory-mediated Overfitting}。这是便于我们讨论的工作术语，并不是该领域已经建立的统一理论名称。

\hypertarget{part-3-34-0001-memory-compression-ux662fux5426ux53efux4ee5ux51cfux5c11ux8fc7ux62dfux5408}{%
\subsubsection{3.5 Memory Compression 是否可以减少过拟合？}\label{part-3-34-0001-memory-compression-ux662fux5426ux53efux4ee5ux51cfux5c11ux8fc7ux62dfux5408}}

假设 Agent 不再保存全部实验反馈，而是将历史记录压缩为一份研究摘要：
\[M_k=\mathcal S(H_k),\]

其中 \(\mathcal S\) 是 Summary Mechanism。

如果 Memory 仅由历史记录生成，不存在其他来自评估集的信息输入，则根据 Data Processing Inequality：
\[\boxed{
I(D_{\mathrm{IS}};M_k)
\leq
I(D_{\mathrm{IS}};H_k)
}\]

这意味着压缩后的 Memory 不会比原始历史记录包含更多关于评估数据的信息。

但这里有一个 Nontrivial Insight：

\textbf{信息量减少，不意味着过拟合一定减少。}

假设原始记录包含大量实验细节，而摘要只保留了：

\begin{quote}
在开发数据中表现最好的特征组合是 \(X_3,X_7,X_{12}\)。
\end{quote}

虽然信息量大幅减少，但最具有选择偏差的信息可能恰好被保留了下来。

因此，我们研究 Memory 时不能只测量其长度，还应当分析其内容、来源和使用方式。

\begin{readertable}
\begin{longtable}[]{@{}
  >{\raggedright\arraybackslash}p{(\columnwidth - 2\tabcolsep) * \real{0.36}}
  >{\raggedright\arraybackslash}p{(\columnwidth - 2\tabcolsep) * \real{0.60}}@{}}
\toprule
Memory Attribute & Potential Research Variable \\
\midrule
\endhead
Capacity & 保存多少历史信息 \\
Granularity & 标量分数、样本级错误或完整日志 \\
Persistence & 信息保留多少轮 \\
Retrieval & 选择哪些历史记录 \\
Compression & 如何提炼经验 \\
Provenance & 信息来自哪个数据集和哪次实验 \\
\bottomrule
\end{longtable}
\end{readertable}

这也为后续 Memory Ablation Study 提供了实验变量。

\hypertarget{part-3-34-0001-planning-agent-ux5982ux4f55ux51b3ux5b9aux4e0bux4e00ux6b65ux884cux52a8}{%
\subsection{4. Planning --- Agent 如何决定下一步行动？}\label{part-3-34-0001-planning-agent-ux5982ux4f55ux51b3ux5b9aux4e0bux4e00ux6b65ux884cux52a8}}

\hypertarget{part-3-34-0001-planning-ux7684ux57faux672cux542bux4e49}{%
\subsubsection{4.1 Planning 的基本含义}\label{part-3-34-0001-planning-ux7684ux57faux672cux542bux4e49}}

Planning 指根据当前目标、状态和可用资源，组织未来的一系列行动。

例如，研究任务为：

\begin{quote}
Improve the OOS robustness of a financial prediction model.
\end{quote}

Agent 可能形成如下计划：

检查现有模型及 Baseline Performance。

分析可能存在的 Feature Redundancy。

构造三个候选正则化方案。

运行 Chronological Validation。

比较候选方案并保存实验结果。

根据反馈决定是否继续探索。

如果整个计划始终由人类预先写死，它更接近 Workflow。如果 Agent 可以根据实验结果调整后续计划，就具有了更强的 Agentic Behavior。

\hypertarget{part-3-34-0001-representative-planning-strategies}{%
\subsubsection{4.2 Representative Planning Strategies}\label{part-3-34-0001-representative-planning-strategies}}

\begin{readertable}
\begin{longtable}[]{@{}
  >{\raggedright\arraybackslash}p{(\columnwidth - 4\tabcolsep) * \real{0.31}}
  >{\raggedright\arraybackslash}p{(\columnwidth - 4\tabcolsep) * \real{0.35}}
  >{\raggedright\arraybackslash}p{(\columnwidth - 4\tabcolsep) * \real{0.31}}@{}}
\toprule
Strategy & Basic Idea & Characteristic \\
\midrule
\endhead
Sequential Planning & 按照预定顺序执行 & 简单、容易复现 \\
Replanning & 根据新观察修改计划 & 具有动态适应性 \\
Tree Search & 同时探索多个候选路径 & 搜索范围较大 \\
Hierarchical Planning & 将任务分解为多个子任务 & 适合复杂研究 \\
ReAct & 交替组织推理、行动和观察 & 根据环境反馈持续调整 \\
\bottomrule
\end{longtable}
\end{readertable}

\hypertarget{part-3-34-0001-react}{%
\paragraph{ReAct}\label{part-3-34-0001-react}}

ReAct: Synergizing Reasoning and Acting in Language Models

是理解 Agent 动态规划的一项经典研究。

其主要思想是将 Reasoning 与 Acting 交替组织，使外部观察能够影响下一步决策。

\href{https://arxiv.org/abs/2210.03629}{来源 1}

\begin{verbatim}
Reasoning / Decision
    ↓
Action
    ↓
Observation
    ↓
Updated Decision
    ↓
Next Action
\end{verbatim}

例如：

\begin{readertable}
\begin{longtable}[]{@{}
  >{\raggedright\arraybackslash}p{(\columnwidth - 2\tabcolsep) * \real{0.31}}
  >{\raggedright\arraybackslash}p{(\columnwidth - 2\tabcolsep) * \real{0.65}}@{}}
\toprule
Step & Agent Behavior \\
\midrule
\endhead
Decision & 怀疑模型表现不佳来自特征共线性 \\
Action & 计算 Correlation Matrix \\
Observation & 发现多个特征高度相关 \\
Decision & 尝试 Ridge Regression \\
Action & 运行比较实验 \\
Observation & 记录风险与评估得分 \\
\bottomrule
\end{longtable}
\end{readertable}

这个过程本身并不要求修改底层模型权重。

\hypertarget{part-3-34-0001-planning-ux4e0e-explorationexploitation-trade-off}{%
\subsubsection{4.3 Planning 与 Exploration--Exploitation Trade-off}\label{part-3-34-0001-planning-ux4e0e-explorationexploitation-trade-off}}

这一部分与你以后可能接触的 Reinforcement Learning 有一定联系，但目前不需要学习完整的 RL 理论。

Agent 在选择下一轮研究方向时，通常面临两种选择：

\begin{readertable}
\begin{longtable}[]{@{}
  >{\raggedright\arraybackslash}p{(\columnwidth - 2\tabcolsep) * \real{0.35}}
  >{\raggedright\arraybackslash}p{(\columnwidth - 2\tabcolsep) * \real{0.61}}@{}}
\toprule
Strategy & Meaning \\
\midrule
\endhead
Exploration & 尝试此前没有研究过的新方法 \\
Exploitation & 继续改进目前表现较好的方法 \\
\bottomrule
\end{longtable}
\end{readertable}

\bookgraphic{0.544}{part-3-34-0001-007-01}{}

对于 RSI，这两种搜索策略都可能带来统计风险。

过度 Exploration 可能引入大量候选模型，增加 Multiple Testing；过度 Exploitation 则可能不断深入某个由历史噪声驱动的错误方向。

所以，简单增加 Exploration 或 Exploitation 都不能保证改善 Generalization。

我们会在 III-D 中系统研究 Search Strategy。

\hypertarget{part-3-34-0001-reflection-agent-ux5982ux4f55ux6839ux636eux53cdux9988ux5b66ux4e60}{%
\subsection{5. Reflection --- Agent 如何根据反馈学习？}\label{part-3-34-0001-reflection-agent-ux5982ux4f55ux6839ux636eux53cdux9988ux5b66ux4e60}}

这是本阶段第二个重点。

\hypertarget{part-3-34-0001-self-reflection-ux7684ux5b9aux4e49}{%
\subsubsection{5.1 Self-Reflection 的定义}\label{part-3-34-0001-self-reflection-ux7684ux5b9aux4e49}}

Self-Reflection 指 Agent 对此前的行动、结果和失败原因进行分析，并尝试形成能够改善后续行为的反馈。

假设 Agent 第一次实验失败：

\begin{verbatim}
Validation MSE increased from 0.42 to 0.57.
\end{verbatim}

Agent 可能生成如下 Reflection：

\begin{verbatim}
The additional features may have increased
estimation variance without introducing
sufficient predictive information.

The next experiment should examine feature
regularization and redundancy.
\end{verbatim}

注意，这份 Reflection 并不是经过证明的实验结论。

它只是 Agent 根据当前证据生成的一个解释，以及一个下一步研究假设。

这是我们分析 Research Agent 时必须保持的认识：

\textbf{Reflection is hypothesis generation, not hypothesis verification.}

\hypertarget{part-3-34-0001-three-sources-of-feedback}{%
\subsubsection{5.2 Three Sources of Feedback}\label{part-3-34-0001-three-sources-of-feedback}}

\begin{readertable}
\begin{longtable}[]{@{}
  >{\raggedright\arraybackslash}p{(\columnwidth - 4\tabcolsep) * \real{0.23}}
  >{\raggedright\arraybackslash}p{(\columnwidth - 4\tabcolsep) * \real{0.33}}
  >{\raggedright\arraybackslash}p{(\columnwidth - 4\tabcolsep) * \real{0.41}}@{}}
\toprule
Feedback Source & Example & Reliability Consideration \\
\midrule
\endhead
Intrinsic Feedback & LLM 自己检查自己的答案 & 可能重复原来的推理错误 \\
External Feedback & Unit Test、程序输出、评估得分 & 可提供客观证据，但可能存在统计噪声 \\
Human Feedback & 研究人员评价实验设计 & 可以引入领域知识，但也可能存在主观偏差 \\
\bottomrule
\end{longtable}
\end{readertable}

其中 External Feedback 并不等于 Ground Truth。

例如，对于一个确定性的算法题，Unit Test 的失败通常可以说明程序存在具体问题。

但对于金融预测任务，某个模型在有限验证集上表现下降，并不必然意味着模型的真实性能下降。

这种区别对于理解金融 RSI 尤其重要。

\hypertarget{part-3-34-0001-self-refine}{%
\subsubsection{5.3 Self-Refine}\label{part-3-34-0001-self-refine}}

2023 年发表于 NeurIPS 的 \emph{Self-Refine: Iterative Refinement with Self-Feedback} 提出了一种利用同一个语言模型生成、评价并修改输出的迭代方法，无需为该流程额外训练模型。

\href{https://proceedings.neurips.cc/paper_files/paper/2023/hash/91edff07232fb1b55a505a9e9f6c0ff3-Abstract-Conference.html}{来源 1}

其基本结构为：

\textbf{Initial Output}

\(y_0\) ↓ \textbf{Self-Feedback}

\(f_k=\operatorname{Feedback}(x,y_k)\) ↓ \textbf{Refinement}

\(y_{k+1}=\operatorname{Refine}(x,y_k,f_k)\)

重复直到满足预先定义的停止条件。

Self-Refine 主要研究的是 Output-level Improvement，即如何改进当前输出。

它并不必然意味着 Agent 的 Prompt、Control Logic 或自我改进机制发生了持久性改变。

\hypertarget{part-3-34-0001-reflexion}{%
\subsubsection{5.4 Reflexion}\label{part-3-34-0001-reflexion}}

与本课题联系更加紧密的是：

\textbf{Reflexion: Language Agents with Verbal Reinforcement Learning}

该研究发表于 NeurIPS 2023，利用语言形式的反馈与 Episodic Memory 改善 Agent 后续尝试的表现，而不是直接通过梯度更新底层模型权重。

\href{https://proceedings.neurips.cc/paper_files/paper/2023/hash/1b44b878bb782e6954cd888628510e90-Abstract-Conference.html}{来源 1} \href{https://arxiv.org/abs/2303.11366}{来源 2}

\textbf{Reflexion --- Simplified Architecture}

\bookgraphic{0.486}{part-3-34-0001-011-01}{根据论文核心思想重新绘制的简化架构，不是原论文的完整实现图。}

Reflexion 的重要性在于，它展示了一种无需更新模型参数的学习机制：
\[\boxed{
\text{Learning through Memory and Feedback}
}\]

论文标题使用了 Verbal Reinforcement Learning，但它与通过 Policy Gradient、Q-learning 等方法更新模型参数的传统 RL 并不相同。

对于 X Program \#161，这提示我们可以研究：

\begin{itemize}
\tightlist
\item
  Agent 如何将实验反馈转化为语言形式的经验；
\item
  这些经验如何影响下一轮候选方法的生成；
\item
  哪些经验具有跨环境的可迁移性；
\item
  哪些经验只是对有限开发数据的过度适应。
\end{itemize}

\hypertarget{part-3-34-0001-self-reflection-ux771fux7684ux80fdux591fux4fddux8bc1ux6539ux5584ux5417}{%
\subsubsection{5.5 Self-Reflection 真的能够保证改善吗？}\label{part-3-34-0001-self-reflection-ux771fux7684ux80fdux591fux4fddux8bc1ux6539ux5584ux5417}}

不能。

2024 年发表于 ICLR 的 \emph{Large Language Models Cannot Self-Correct Reasoning Yet} 研究了 Intrinsic Self-Correction，即 LLM 在没有外部反馈时修正自身推理的能力。论文报告，在所研究的推理任务中，这种自我纠正存在明显困难，某些情况下甚至导致性能下降。

\href{https://proceedings.iclr.cc/paper_files/paper/2024/hash/8b4add8b0aa8749d80a34ca5d941c355-Abstract-Conference.html}{来源 1} \href{https://openreview.net/forum?id=IkmD3fKBPQ}{来源 2}

这与 Self-Refine 的正面实验结果并不必然矛盾，因为二者研究的任务、反馈方式和实验协议并不完全相同。

我们应该保留的结论是：
\[\boxed{
\text{More Reflection}
\not\Rightarrow
\text{More Reliable Improvement}
}\]

Reflection 的有效性取决于反馈质量、任务性质、评价机制以及后续修改是否经过可靠验证。

\hypertarget{part-3-34-0001-self-correctionself-reflection-ux4e0e-self-modification}{%
\subsection{6. Self-Correction、Self-Reflection 与 Self-Modification}\label{part-3-34-0001-self-correctionself-reflection-ux4e0e-self-modification}}

这三个术语非常容易混淆。进入 III-C 之前，我们需要将它们明确区分。

\begin{readertable}
\begin{longtable}[]{@{}
  >{\raggedright\arraybackslash}p{(\columnwidth - 6\tabcolsep) * \real{0.27}}
  >{\raggedright\arraybackslash}p{(\columnwidth - 6\tabcolsep) * \real{0.27}}
  >{\raggedright\arraybackslash}p{(\columnwidth - 6\tabcolsep) * \real{0.15}}
  >{\raggedright\arraybackslash}p{(\columnwidth - 6\tabcolsep) * \real{0.27}}@{}}
\toprule
Concept & 改变的对象 & Persistent? & Example \\
\midrule
\endhead
Self-Correction & 当前输出 & 不一定 & 修复当前 Python 代码 \\
Self-Reflection & 对过去行动的解释 & 不一定 & 分析实验失败原因 \\
Memory-based Learning & 历史经验与行为上下文 & 可以 & 保存研究经验供下一次使用 \\
Self-Modification & Agent 自身的方法或结构 & 通常是 & 修改永久性的搜索流程 \\
Recursive Self-Improvement & 能够继续参与自我改进的系统 & 是 & 改进后的优化器继续改进自身 \\
\bottomrule
\end{longtable}
\end{readertable}

\textbf{Output-level Change}

修改当前答案或解决方案

↓ \textbf{Behavior-level Change}

通过反馈和 Memory 改变未来行为

↓ \textbf{Method-level Change}

修改持久性的 Prompt、程序或工作机制

↓ \textbf{Recursive Improvement}

改进后的方法继续参与自身的改进

这是一种用于理解系统层级的概念框架，并不是所有 Agent 必须依次经历的四个发展阶段。

尤其重要的是：
\[\text{Self-Reflection}
\not\Rightarrow
\text{Self-Modification}.\]

一个系统可以不断反思，却始终没有修改自己的持久方法。即使发生了 Self-Modification，也不意味着每一次修改都是真正的 Improvement。

\hypertarget{part-3-34-0001-ux4e00ux4e2aux5b8cux6574ux7684-financial-rsi-failure-example}{%
\subsection{7. 一个完整的 Financial RSI Failure Example}\label{part-3-34-0001-ux4e00ux4e2aux5b8cux6574ux7684-financial-rsi-failure-example}}

现在尝试将本阶段的四个组件与 Part II 的统计学知识连接起来。

假设我们拥有一个 Financial Research Agent，其初始任务是寻找能够预测未来收益率的特征。

\hypertarget{part-3-34-0001-agent-ux7684ux6f14ux5316ux8fc7ux7a0b}{%
\subsubsection{7.1 Agent 的演化过程}\label{part-3-34-0001-agent-ux7684ux6f14ux5316ux8fc7ux7a0b}}

\textbf{Round 0 --- Baseline}

\begin{verbatim}
  使用少量具有明确经济含义的特征，获得初始性能。
\end{verbatim}

\textbf{Round 1 --- Tool Expansion}

\begin{verbatim}
  新增自动特征搜索工具，发现若干真正具有预测价值的变量。
\end{verbatim}

\textbf{Round 2 --- Planning Improvement}

\begin{verbatim}
  扩大候选搜索范围，并开始比较不同窗口长度及特征组合。
\end{verbatim}

\textbf{Round 3 --- Reflection}

\begin{verbatim}
  Agent 将偶然获得的 IS Improvement 解释为增加特征交互项的有效性。
\end{verbatim}

\textbf{Round 4 --- Memory Update}

\begin{verbatim}
  将这一未经充分验证的经验保存为长期研究原则。
\end{verbatim}

\textbf{Round 5 --- Self-Modification}

\begin{verbatim}
  修改特征生成程序，使未来的搜索优先考虑这类交互项。
\end{verbatim}

\textbf{Round 6+ --- Reinforcement of Search Bias}

\begin{verbatim}
  后续实验继续围绕相似的历史规律进行优化，可能进一步损害 OOS Performance。
\end{verbatim}

\bookgraphic{0.786}{part-3-34-0001-015-01}{}

假想实验结果，用于说明 Tools、Planning、Reflection、Memory 和 Self-Modification 可能如何共同形成 Better, Then Worse 的演化轨迹。

\hypertarget{part-3-34-0001-ux5bf9ux5e94ux7684ux7edfux8ba1ux673aux5236}{%
\subsubsection{7.2 对应的统计机制}\label{part-3-34-0001-ux5bf9ux5e94ux7684ux7edfux8ba1ux673aux5236}}

\begin{readertable}
\begin{longtable}[]{@{}
  >{\raggedright\arraybackslash}p{(\columnwidth - 4\tabcolsep) * \real{0.25}}
  >{\raggedright\arraybackslash}p{(\columnwidth - 4\tabcolsep) * \real{0.32}}
  >{\raggedright\arraybackslash}p{(\columnwidth - 4\tabcolsep) * \real{0.38}}@{}}
\toprule
Agent Component & System-level Change & Statistical Mechanism \\
\midrule
\endhead
Tools & 扩大候选特征集合 & Hypothesis Space Expansion \\
Planning & 根据历史结果调整搜索方向 & Adaptive Search \\
Reflection & 将偶然成功解释为有效规律 & Selection Bias Misinterpretation \\
Memory & 保存并重复利用这一解释 & Information Accumulation \\
Self-Modification & 将解释写入持久性程序 & Persistent Search Bias \\
\bottomrule
\end{longtable}
\end{readertable}

这与我们在 Part II 中所建立的理论是一致的。

但这里增加了一个新的认识：\textbf{统计偏差可以通过软件系统中的不同组件传播，并最终转化为持久性的行为偏差。}

一次错误判断并不可怕，真正值得研究的是 Agent 是否会通过自身的反馈机制反复强化这次错误判断。

以上是待验证的机制假设，而非已经观察到的课题组实验结果。

\hypertarget{part-3-34-0001-research-practice-memory-reflection-ablation}{%
\subsection{8. Research Practice --- Memory \& Reflection Ablation}\label{part-3-34-0001-research-practice-memory-reflection-ablation}}

如果我们希望研究 Memory 和 Reflection 是否加重 RSI Overfitting，可以设计一个 \(2\times2\) Factorial Experiment。

\hypertarget{part-3-34-0001-experimental-groups}{%
\subsubsection{8.1 Experimental Groups}\label{part-3-34-0001-experimental-groups}}

\textbf{Group A --- Baseline}

Persistent Memory

Reflection

\textbf{Group B --- Memory Only}

Persistent Memory

Reflection

\textbf{Group C --- Reflection Only}

Persistent Memory

Reflection

\textbf{Group D --- Full System}

Persistent Memory

Reflection

这里关闭 Persistent Memory，并不意味着完全删除 Agent 的当前状态或代码版本，而是禁止跨轮次检索额外的实验历史和经验记录。

关闭 Reflection 则意味着不生成额外的自然语言实验分析，Agent 仍然可以接收必要的数值评分。

尽量保持各组底层模型、初始 Agent、任务、评估协议和搜索预算相同，同时记录实际 Token 和计算开销。

\hypertarget{part-3-34-0001-main-effect-ux4e0e-interaction-effect}{%
\subsubsection{8.2 Main Effect 与 Interaction Effect}\label{part-3-34-0001-main-effect-ux4e0e-interaction-effect}}

这实际上是你可以利用已有 Applied Statistics 基础处理的问题。

设 \(Y\) 为某个实验结果，例如 Final OOS Score。

定义：
\[x_M=
\begin{cases}
1,&\text{Memory enabled},\\
0,&\text{Memory disabled},
\end{cases}\]

\[x_R=
\begin{cases}
1,&\text{Reflection enabled},\\
0,&\text{Reflection disabled}.
\end{cases}\]

我们可以建立一个简单的 Factorial Regression Model：
\[\boxed{
Y
=
\beta_0
+\beta_Mx_M
+\beta_Rx_R
+\beta_{MR}x_Mx_R
+\varepsilon
}\]

其中：

\begin{readertable}
\begin{longtable}[]{@{}
  >{\raggedright\arraybackslash}p{(\columnwidth - 2\tabcolsep) * \real{0.32}}
  >{\raggedright\arraybackslash}p{(\columnwidth - 2\tabcolsep) * \real{0.64}}@{}}
\toprule
Coefficient & Interpretation \\
\midrule
\endhead
\(\beta_0\) & Baseline Performance \\
\(\beta_M\) & 无 Reflection 时 Memory 的影响 \\
\(\beta_R\) & 无 Memory 时 Reflection 的影响 \\
\(\beta_{MR}\) & Memory 与 Reflection 的交互效应 \\
\bottomrule
\end{longtable}
\end{readertable}

假设四组总体平均 OOS 表现分别为 \(\mu_A,\mu_B,\mu_C,\mu_D\)，那么：
\[\boxed{
\beta_{MR}
=
\mu_D-\mu_B-\mu_C+\mu_A
}\]

这个 Interaction Effect 尤其值得研究。

因为我们可能发现，Memory 或 Reflection 单独使用时具有一定好处，但二者结合后，却通过反复积累和强化错误经验导致后期性能下降。

\textbf{Factorial Experiment --- Interaction Effect}

输入四组平均 OOS Score，观察 Memory 和 Reflection 的交互关系。初始值为假想数据。

A --- Baseline 0.60

B --- Memory 0.67

C --- Reflection 0.65

D --- Memory + Reflection 0.56

\bookgraphic{0.531}{part-3-34-0001-018-01}{}

Without Reflection

With Reflection

βₘ

\textbf{0.070}

βᵣ

\textbf{0.050}

βₘᵣ

\textbf{-0.160}

非平行的两条线表示存在观察到的交互效应，但是否具有统计显著性仍需要重复实验及误差分析。

我认为这是一个比较适合本科生参与的研究切入点。因为它不要求你直接设计新的 LLM Architecture，而是要求你理解 Agent 组件、控制实验变量，并对结果进行统计分析。

当然，真实实验不能只有四个观测值。每一组都需要多个独立 RSI Runs，并且需要考虑共同随机种子、市场窗口和其他依赖结构。

\hypertarget{part-3-34-0001-literature-index}{%
\subsection{9. Literature Index}\label{part-3-34-0001-literature-index}}

本阶段涉及四篇已经发表的同行评审论文。

\textbf{P0 --- Reflexion: Language Agents with Verbal Reinforcement Learning}

Noah Shinn et al. · NeurIPS 2023

本阶段优先级最高的论文。重点理解 Actor、Evaluator、Self-Reflection 与 Episodic Memory 之间的关系。

\href{https://arxiv.org/abs/2303.11366}{arXiv:2303.11366} · \href{https://proceedings.neurips.cc/paper_files/paper/2023/hash/1b44b878bb782e6954cd888628510e90-Abstract-Conference.html}{正式出版页面}

\textbf{P1 --- ReAct: Synergizing Reasoning and Acting in Language Models}

Shunyu Yao et al. · ICLR 2023

理解 Reasoning--Action--Observation Loop。

\href{https://arxiv.org/abs/2210.03629}{arXiv:2210.03629} · \href{https://arxiv.org/abs/2210.03629}{来源 1}

\textbf{P1 --- Self-Refine: Iterative Refinement with Self-Feedback}

Aman Madaan et al. · NeurIPS 2023

理解 Output-level Self-Improvement，以及如何使用语言形式的反馈。

\href{https://arxiv.org/abs/2303.17651}{arXiv:2303.17651} · \href{https://proceedings.neurips.cc/paper_files/paper/2023/hash/91edff07232fb1b55a505a9e9f6c0ff3-Abstract-Conference.html}{正式出版页面}

\textbf{P1 --- Large Language Models Cannot Self-Correct Reasoning Yet}

Jie Huang et al. · ICLR 2024

理解 Intrinsic Self-Correction 的局限，避免将自我反思误认为可靠的验证机制。

\href{https://arxiv.org/abs/2310.01798}{arXiv:2310.01798} · \href{https://openreview.net/forum?id=IkmD3fKBPQ}{OpenReview}

建议优先阅读 Reflexion，再阅读第四篇。前者展示反馈与 Memory 如何帮助 Agent 学习，后者则提供了理解自我纠正局限性的另一种视角。

\hypertarget{part-3-34-0001-iii-b-conclusion-index}{%
\subsection{10. III-B --- Conclusion Index}\label{part-3-34-0001-iii-b-conclusion-index}}

\begin{readertable}
\begin{longtable}[]{@{}
  >{\raggedright\arraybackslash}p{(\columnwidth - 4\tabcolsep) * \real{0.13}}
  >{\raggedright\arraybackslash}p{(\columnwidth - 4\tabcolsep) * \real{0.70}}
  >{\raggedright\arraybackslash}p{(\columnwidth - 4\tabcolsep) * \real{0.13}}@{}}
\toprule
ID & Key Conclusion & Index \\
\midrule
\endhead
B1 & Tool Calling 将语言生成转化为实际程序操作 & §2 \\
B2 & Tool Expansion 可能同时扩大统计搜索空间 & §2.3 \\
B3 & Memory 不等于 Context Window & §3.1 \\
B4 & 经验可以逐步转化为持久性的研究方法 & §3.2 \\
B5 & Memory 是历史评估信息影响未来行为的通道 & §3.3 \\
B6 & Memory Compression 不保证缓解过拟合 & §3.5 \\
B7 & Exploration 与 Exploitation 均可能产生统计风险 & §4.3 \\
B8 & Reflection 是生成解释与假设，而不是验证假设 & §5 \\
B9 & Self-Reflection 不等于 Self-Modification & §6 \\
B10 & Factorial Experiment 可以分析 Agent 组件之间的交互效应 & §8 \\
\bottomrule
\end{longtable}
\end{readertable}

\hypertarget{part-3-34-0001-core-formula-sheet}{%
\subsubsection{Core Formula Sheet}\label{part-3-34-0001-core-formula-sheet}}

本阶段不需要掌握新的复杂数学定理，只需保留以下三个数学关系。

\textbf{Result 1 --- Memory Updating}

\[M_{k+1}=U_M(M_k,F_k,H_k).\]

\textbf{Result 2 --- Adaptive Information Flow}

\[D_{\mathrm{IS}}
\longrightarrow F_k
\longrightarrow M_{k+1}
\longrightarrow A_{k+1}.\]

\textbf{Result 3 --- Factorial Experiment}

\[Y=\beta_0+\beta_Mx_M+\beta_Rx_R
+\beta_{MR}x_Mx_R+\varepsilon.\]

\hypertarget{part-3-34-0001-research-takeaway}{%
\subsubsection{Research Takeaway}\label{part-3-34-0001-research-takeaway}}

本阶段最核心的观点，是 Agent 的学习过程可以具有以下两种截然不同的结果：

\textbf{Genuine Learning}

Reliable Feedback

↓

Valid Reflection

↓

Transferable Memory

↓

Generalization Gain

\textbf{Spurious Learning}

Noise-driven Feedback

↓

Incorrect Reflection

↓

Biased Memory

↓

Potential OOS Decay

这两条路径可能具有完全相同的系统架构。

区别不在于 Agent 有没有 Memory 或 Reflection，而在于它所积累的经验是否具有真实的统计有效性。

\hypertarget{part-3-34-0001-next-iii-c-self-modification-recursive-optimization}{%
\subsection{Next --- III-C: Self-Modification \& Recursive Optimization}\label{part-3-34-0001-next-iii-c-self-modification-recursive-optimization}}

III-A --- Agent Architecture \& Harness

III-B --- Tools, Memory, Planning \& Reflection

III-C --- Self-Modification \& Recursive Optimization

III-D --- Evaluation, Search \& Selection

III-E --- Automated Research Agents

III-F --- Failure Modes \& Instrumentation

下一阶段，我们终于开始深入 RSI 本身。

\begin{readertable}
\begin{longtable}[]{@{}
  >{\raggedright\arraybackslash}p{(\columnwidth - 2\tabcolsep) * \real{0.38}}
  >{\raggedright\arraybackslash}p{(\columnwidth - 2\tabcolsep) * \real{0.58}}@{}}
\toprule
Topic & Core Question \\
\midrule
\endhead
Self-Modification & Agent 如何修改自身代码？ \\
Object-level Optimization & 如何改进研究对象？ \\
Meta-level Optimization & 如何改进负责优化的程序？ \\
Recursive Optimization & 如何使用改进后的优化器改进自身？ \\
Code-based Evolution & 如何生成、测试及保留新的 Agent？ \\
STOP & 如何实际实现代码层面的 RSI？ \\
Path Dependence & 为什么不同修改历史可能产生不同的最终结果？ \\
\bottomrule
\end{longtable}
\end{readertable}

届时我们将通过 STOP 的研究思想和简化 Python 伪代码，真正理解：
\[A_{k+1}
=
\mathcal I_{A_k}(A_k,F_k,H_k)\]

这条我们从 Part I 开始就一直使用的抽象公式，究竟如何转化为一个实际可运行的计算系统。
